\chapter{gert Integration}
\label{ch:gert-integration}

This chapter describes how \texttt{gert.home} integrates with the gert runtime: how
property files are parsed, validated, and lowered to core primitives; how away mode is
enforced; how cadence checking works; and what trace events are emitted during execution.

% ─────────────────────────────────────────────────────────────────────────────
\section{Core Integration Hooks}
\label{sec:integration-hooks}
% ─────────────────────────────────────────────────────────────────────────────

A \texttt{.home.yaml} property file passes through four stages before any step executes:

\begin{enumerate}
  \item \textbf{Parse} --- The YAML file is parsed and the \texttt{apiVersion} and
        \texttt{kit} fields are read. The \texttt{home/v0} Kit compiler is activated.
  \item \textbf{Validate} --- The property file is validated against the \texttt{gert.home}
        JSON Schema. Zone and asset references are resolved; missing zone IDs or invalid
        cadence formats produce compile errors.
  \item \textbf{Lower} --- \texttt{home.*} step types are lowered to core gert primitives:
        \texttt{home.routine} → \texttt{iterate} + \texttt{manual}; incident
        \texttt{human\_task} → \texttt{manual}; \texttt{decision} → \texttt{branch};
        \texttt{parallel} → \texttt{parallel}.
  \item \textbf{Execute} --- The gert runtime executes the lowered runbook, checking
        cadence due dates, calling \texttt{IsAwayModeActive()}, and emitting trace events at
        each step boundary.
\end{enumerate}

The gert runtime never sees \texttt{home.*} vocabulary after the lower stage. This means
all core gert tooling (trace readers, evidence bundlers, CI integrations) works with
\texttt{gert.home} runs without modification.

% ─────────────────────────────────────────────────────────────────────────────
\section{Role-Based Gating}
\label{sec:integration-role-gating}
% ─────────────────────────────────────────────────────────────────────────────

\subsection{Executor Role Hints}

The \texttt{executor.role} field on a \texttt{home.routine} is advisory in
\texttt{home/v0}. The runtime reads the field and includes it in notification messages and
the task UI, but does not block execution if the current actor does not match the hint.

\begin{tcolorbox}[colback=yellow!10,colframe=orange!70,title=Note]
  \texttt{home/v0} does not enforce executor role hints at runtime. A homeowner or active
  delegate can always execute any step regardless of the \texttt{executor.role} value. Hard
  role gates are planned for \texttt{home/v1}.
\end{tcolorbox}

\subsection{Delegation Scope Enforcement}

The \texttt{assigns} block in a \texttt{delegations} entry is enforced at runtime.
When \texttt{IsAwayModeActive()} returns \texttt{true}:

\begin{itemize}
  \item The runtime checks whether the current routine's ID appears in any active
        delegation's \texttt{assigns} list as \texttt{\{routine: id\}}.
  \item The runtime also checks whether the routine's \texttt{zone} ID appears as
        \texttt{\{zone: id\}}.
  \item If neither match, the routine is not delegated and the homeowner remains executor.
  \item If a match is found, the delegate's name and contact are substituted as executor
        identity for that run.
\end{itemize}

% ─────────────────────────────────────────────────────────────────────────────
\section{Command / Action Scope}
\label{sec:integration-commands}
% ─────────────────────────────────────────────────────────────────────────────

\texttt{gert.home} routines and incidents generate \texttt{manual} steps --- human-action
steps that pause execution until a person confirms completion and submits evidence.
The kit does not generate \texttt{cli} steps. There are no shell commands, no environment
variable access, and no network calls in a \texttt{gert.home} execution. This keeps the
security surface minimal for household contexts.

% ─────────────────────────────────────────────────────────────────────────────
\section{Away Mode Enforcement}
\label{sec:integration-away-mode}
% ─────────────────────────────────────────────────────────────────────────────

Away mode is evaluated at the start of each routine execution via the
\texttt{IsAwayModeActive()} runtime function:

\begin{enumerate}
  \item The runtime reads all \texttt{delegations} entries from the property file.
  \item It compares the current timestamp against each entry's \texttt{active.from} and
        \texttt{active.to} fields using the timezone from \texttt{GERT\_HOME\_TZ}.
  \item If one or more entries are active, \texttt{IsAwayModeActive()} returns \texttt{true}
        and the runtime proceeds with delegation scope checks.
  \item If no entries are active, \texttt{IsAwayModeActive()} returns \texttt{false} and
        the homeowner is executor for all routines.
\end{enumerate}

When away mode is active, a \texttt{delegation\_active} trace event is emitted for each
routine execution that matches an active delegation scope.

% ─────────────────────────────────────────────────────────────────────────────
\section{Cadence Enforcement}
\label{sec:integration-cadence}
% ─────────────────────────────────────────────────────────────────────────────

Before firing a \texttt{home.routine} step, the runtime checks whether the routine is due:

\begin{enumerate}
  \item The runtime reads the last completion timestamp from the run directory
        (\texttt{runs/\{routine-id\}/latest/trace.jsonl}).
  \item It computes the next due date: last completion + cadence interval (or the seasonal
        interval for the current date).
  \item If the current timestamp is before the due date, the routine is skipped and a
        \texttt{routine\_skipped} trace event is emitted.
  \item If the current timestamp is within \texttt{notifications.remind\_hours\_before} of
        the due date, a reminder notification is sent to the executor.
  \item If the current timestamp is past \texttt{due date + notifications.escalate\_hours\_overdue},
        an overdue alert is sent to the homeowner regardless of delegation state.
\end{enumerate}

% ─────────────────────────────────────────────────────────────────────────────
\section{The Audit Trail}
\label{sec:integration-audit-trail}
% ─────────────────────────────────────────────────────────────────────────────

The \texttt{gert.home} runtime emits structured trace events throughout execution. These are
written to \texttt{runs/\{routine-id\}/\{date\}/trace.jsonl} in newline-delimited JSON.

\begin{description}
  \item[\texttt{routine\_due}] Emitted when a routine's due date arrives. Fields:
        \texttt{routine\_id}, \texttt{zone}, \texttt{executor}.
  \item[\texttt{routine\_completed}] Emitted when all steps of a routine run are completed.
        Fields: \texttt{routine\_id}, \texttt{duration\_minutes}.
  \item[\texttt{incident\_opened}] Emitted when an incident template is instantiated.
        Fields: \texttt{incident\_id}, \texttt{opened\_by}.
  \item[\texttt{incident\_step\_completed}] Emitted when each incident step finishes.
        Fields: \texttt{incident\_id}, \texttt{step\_id}, \texttt{type}.
  \item[\texttt{evidence\_submitted}] Emitted when evidence is captured. Fields:
        \texttt{step\_id}, \texttt{evidence\_type}, \texttt{submitted\_by}, \texttt{sha256}.
  \item[\texttt{delegation\_active}] Emitted at the start of any routine execution where
        an active delegation scope matches. Fields: \texttt{delegate\_name},
        \texttt{routine\_id}, \texttt{active\_until}.
\end{description}

% ─────────────────────────────────────────────────────────────────────────────
\section{Complete Integration Example}
\label{sec:integration-complete-example}
% ─────────────────────────────────────────────────────────────────────────────

The following shows a property file with an active delegation alongside the expected trace
output for a pool-cleaning routine run during the away window.

\begin{minted}{yaml}
apiVersion: runbook/v1
kit: home/v0

property:
  name: Casa Santiago
  zones:
    - id: pool
      label: Swimming Pool

routines:
  - id: pool_clean
    label: Pool cleaning
    zone: pool
    cadence:
      every: 7d
    evidence:
      type: photo
      prompt: Take photo of clear water and chemical test strip
    notifications:
      remind_hours_before: 4
      escalate_hours_overdue: 24

delegations:
  - delegate:
      name: Carlos Méndez
      contact: "+56 9 8765 4321"
    active:
      from: 2025-04-25
      to: 2025-05-02
    assigns:
      - zone: pool
    permissions:
      can_report_incidents: true
      can_modify_routines: false
      can_view_history: true
\end{minted}

Trace events emitted on 2025-04-28 when the routine fires:

\begin{minted}{yaml}
{"seq":1,"ts":"2025-04-28T08:00:00Z","type":"delegation_active",
 "delegate_name":"Carlos Méndez","routine_id":"pool_clean",
 "active_until":"2025-05-02"}
{"seq":2,"ts":"2025-04-28T08:00:01Z","type":"routine_due",
 "routine_id":"pool_clean","zone":"pool","executor":"Carlos Méndez"}
{"seq":3,"ts":"2025-04-28T09:15:00Z","type":"evidence_submitted",
 "step_id":"pool_clean","evidence_type":"photo",
 "submitted_by":"Carlos Méndez",
 "sha256":"a3f2...c1d9"}
{"seq":4,"ts":"2025-04-28T09:15:01Z","type":"routine_completed",
 "routine_id":"pool_clean","duration_minutes":75}
\end{minted}
